%=============================================================================
% WAVE 3, ITERATION 7: PATENT 1 + PATENT 2 + PATENT 11 COMBINED UPDATE
% P1: Field Drag / Gravitational Cloaking --- DEFINITIVE FINAL ASSESSMENT
% P2: Resonators and Metamaterials --- LIGO Systematic Budget
% P11: EM Coupling / SRF Cavity --- SQMS Phase I Protocol, Q_psi Constraints
%
% Agent 1/2/11 | DFD Research Programme
% Building on: W3i6_P1_P2_update.tex, W3i6_P8_P11_update.tex,
%              W3i6_MatSci_update.tex
% Date: 20 March 2026
%
% W3i7 MANDATE:
%   1. LIGO 6th Channel: COMPLETE systematic error budget. Dominant
%      systematics, distinguishability from true psi signal, null tests,
%      realistic reach vs systematic floor.
%   2. SQMS Phase I: Detailed experimental protocol. Measurements,
%      analysis, detection vs null criteria.
%   3. Indirect Q_psi measurement: Can SQMS multi-frequency Q0 data
%      already constrain Q_psi? Existing datasets for reanalysis.
%   4. P1 FINAL STATUS: Definitive assessment with all corrections
%      propagated.
%
% Status flags: [VERIFIED], [CONJECTURE], [INVALIDATED], [CORRECTED], [NEW]
%=============================================================================

\documentclass[11pt]{article}
\usepackage[utf8]{inputenc}
\usepackage{amsmath,amssymb,amsfonts,amsthm,mathtools}
\usepackage{geometry}
\geometry{margin=1in}
\usepackage{booktabs}
\usepackage{siunitx}
\usepackage{hyperref}
\usepackage{xcolor}
\usepackage{enumitem}
\usepackage{float}
\usepackage{array}
\usepackage{longtable}
\usepackage{tcolorbox}
\tcbuselibrary{skins,breakable}

\newcommand{\dpsi}{\delta\psi}
\newcommand{\lam}{\lambda}
\newcommand{\Opsi}{\Omega_\psi}
\newcommand{\gpsi}{\gamma_\psi}
\newcommand{\Mpsi}{M_\psi}
\newcommand{\Meff}{M_{\rm eff}}
\newcommand{\order}[1]{\mathcal{O}(#1)}
\newcommand{\ee}[1]{\times 10^{#1}}
\newcommand{\half}{\tfrac{1}{2}}
\newcommand{\kB}{k_{\mathrm{B}}}
\newcommand{\Qpsi}{Q_\psi}
\newcommand{\SNR}{\mathrm{SNR}}
\newcommand{\Heff}{H_n}
\newcommand{\Ucav}{U_0}
\newcommand{\Vcav}{V_{\rm cav}}
\newcommand{\muG}{\mu_0^{\rm gal}}
\newcommand{\lbone}{|\lam - 1|}
\newcommand{\psicosmo}{\bar{\psi}_{\rm cosmo}}
\newcommand{\psiEM}{\bar{\psi}_{\rm EM}}
\newcommand{\psigrav}{\psi_{\rm grav}}
\newcommand{\dpsiEM}{\delta\psi_{\rm EM}}
\newcommand{\GN}{G_N}

\definecolor{strongcolor}{rgb}{0,0.45,0}
\definecolor{weakcolor}{rgb}{0.65,0,0}
\definecolor{conjcolor}{rgb}{0.6,0.3,0}
\definecolor{rowhi}{rgb}{0.85,0.95,0.85}
\definecolor{rowmed}{rgb}{0.95,0.95,0.80}
\definecolor{rowlo}{rgb}{0.97,0.87,0.87}
\definecolor{falsified}{rgb}{0.7,0,0}

\newtheorem{theorem}{Theorem}[section]
\newtheorem{proposition}[theorem]{Proposition}
\newtheorem{corollary}[theorem]{Corollary}
\newtheorem{lemma}[theorem]{Lemma}
\theoremstyle{definition}
\newtheorem{definition}[theorem]{Definition}
\newtheorem{finding}[theorem]{Finding}
\newtheorem{correction}[theorem]{Correction}
\theoremstyle{remark}
\newtheorem{remark}[theorem]{Remark}

\title{Wave 3, Iteration 7: Patent 1 (Field Drag) + Patent 2 (Resonators)
+ Patent 11 (EM Coupling)\\
\large LIGO 6th-Channel Complete Systematic Budget,\\
SQMS Phase~I Experimental Protocol,\\
Indirect $\Qpsi$ Constraints from Existing Data,\\
P1 Definitive Final Assessment}

\author{DFD Research Programme --- Agent 1/2/11 (W3i7)\\
\textit{Building on: W3i6 P1/P2, W3i6 P8/P11, W3i6 MatSci}}

\date{20 March 2026}

\begin{document}
\maketitle
\tableofcontents
\newpage

%=============================================================================
\section*{W3i7 Context: State Inherited from W3i6}
%=============================================================================

\noindent\textbf{Critical results from W3i6 carried forward:}

\begin{enumerate}[nosep]
\item \textbf{Passive $\gg$ active by 14 orders of magnitude.}
  LIGO 6th channel: $\lbone \sim 6\ee{-19}$; best active (MEMS+P11):
  $\lbone \sim 10^{-5}$.
\item \textbf{P11 three-phase detection roadmap.}
  Phase~I: SQMS 4-cavity (\$2--4M, 2027--28, $7.8\ee{-10}$);
  Phase~II: MEMS hybrid (2028--30, $4.7\ee{-11}$);
  Phase~III: 256-cavity or ng MEMS (2030--35, $\sim 10^{-12}$).
\item \textbf{$\Qpsi \geq 4\ee{10}$ required} for MEMS+P11 with
  1~$\mu$m NEMS resonator.
\item \textbf{Passive Nb dominates} the passive material FOM
  (inverts the active YBCO ranking).
\item \textbf{P1 experimentally dead.} Signal $\sim 10^{-27}$~s,
  $10^{15}\times$ below detection. Theoretically self-consistent
  after two-component T5 resolution.
\item \textbf{Five-step passive detection programme:}
  (1) LIGO archival, (2) SQMS multi-frequency, (3) ESS Lund,
  (4) indirect $\Qpsi$, (5) conditional MEMS+P11.
\end{enumerate}

\noindent\textbf{W3i7 mandate}: Four deep dives that W3i6 flagged as
open questions.

\begin{tcolorbox}[colback=blue!5!white, colframe=blue!60!black,
  title=\textbf{W3i7 MANDATE}]
\begin{enumerate}[nosep,label=\textbf{M\arabic*.}]
\item \textbf{LIGO 6th channel}: Derive the \emph{complete} systematic
  error budget. What are the dominant systematics? Can they be
  distinguished from a true psi signal? What null tests exist?
  Is $6\ee{-19}$ realistic or does a systematic floor limit reach?
\item \textbf{SQMS Phase~I}: Detailed experimental protocol. What
  measurements, what analysis, what constitutes detection vs null?
\item \textbf{Indirect $\Qpsi$ measurement}: Can SQMS multi-frequency
  $Q_0$ data already constrain $\Qpsi$? What existing datasets should
  be reanalysed?
\item \textbf{P1 FINAL STATUS}: With all corrections propagated, write
  the definitive P1 assessment. What survives, what is dead, what is
  theoretical only?
\end{enumerate}
\end{tcolorbox}


%=============================================================================
%=============================================================================
\part{LIGO 6th Channel: Complete Systematic Error Budget}
\label{part:LIGO}
%=============================================================================
%=============================================================================

%=============================================================================
\section{The 6th-Channel Signal: Precise Prediction}
\label{sec:ligo_signal}
%=============================================================================

\subsection{Signal origin in the DFD framework}

The ``6th channel'' refers to a scalar breathing mode that DFD predicts
in addition to the two GR tensor polarisations ($h_+, h_\times$) and the
three additional polarisations permitted by generic metric theories (vector
$h_x, h_y$ and scalar longitudinal $h_b$). In DFD, the scalar $\psi$-field
perturbation modifies the effective metric:
\begin{equation}
ds^2 = -e^{2\psi}\,c^2 dt^2 + e^{-2\psi}\,dx^2,
\label{eq:DFD_metric}
\end{equation}
so a propagating $\psi$-perturbation $\delta\psi(t,\mathbf{x})$ produces
an isotropic strain:
\begin{equation}
h_\psi = 2\,\delta\psi.
\label{eq:h_psi}
\end{equation}

In a Michelson interferometer with perpendicular arms, the tensor GW mode
$h_+$ produces a differential arm-length change $\Delta L = h_+ L/2$.
The scalar breathing mode produces an \emph{equal} arm-length change in
both arms:
\begin{equation}
\Delta L_x = \Delta L_y = h_\psi\,L = 2\,\delta\psi\,L.
\label{eq:common_mode}
\end{equation}

\begin{theorem}[LIGO Scalar Channel Sensitivity]
\label{thm:ligo_scalar}
A power-recycled Fabry--Perot Michelson interferometer (LIGO
configuration) is \emph{insensitive} to a pure scalar breathing mode
in its primary differential channel. The psi-signal appears only
through:
\begin{enumerate}[nosep]
\item \textbf{Arm asymmetry}: If the two arms have length difference
  $\Delta L_{\rm asym}$, the scalar mode leaks into the differential
  output with suppression factor $\Delta L_{\rm asym}/L$.
  For LIGO: $\Delta L_{\rm asym}/L \lesssim 10^{-3}$
  (actively controlled).
\item \textbf{Frequency-dependent coupling}: If $\delta\psi$ has
  spatial gradients across the detector (wavelength $\lambda_\psi
  \lesssim L$), the scalar response develops an arm-dependent phase:
  \begin{equation}
  h_{\rm diff}^{\rm scalar} = 2\,\delta\psi\,\sin^2\!\left(
  \frac{\pi L}{\lambda_\psi}\right)\,\cos\theta,
  \end{equation}
  where $\theta$ is the angle between $\hat{k}_\psi$ and the
  arm bisector.
\item \textbf{Process~A coupling in mirror coatings}: EM energy stored
  in the mirror coatings experiences a $(\lam-1)$-dependent refractive
  index shift $\delta n = (\lam-1)\,\delta\psi$, producing an optical
  path phase shift $\delta\phi = (2\pi/\lambda_{\rm laser})\cdot
  (\lam-1)\,\delta\psi\cdot d_{\rm coat}$, where $d_{\rm coat}
  \approx 5~\mu$m is the coating thickness.
\end{enumerate}
\end{theorem}

\begin{finding}[Revised LIGO 6th-Channel Coupling]
\label{find:ligo_revised}
The W3i5/W3i6 estimate of $\lbone \sim 6\ee{-19}$ assumed that the
full LIGO strain sensitivity translates directly to a psi-constraint.
This is \textbf{incorrect}. The actual coupling is suppressed by:

\textbf{Channel 1 (arm asymmetry)}: $h_{\rm eff} = 2\delta\psi\cdot
(\Delta L_{\rm asym}/L) \sim 2\delta\psi\times 10^{-3}$. The bound
becomes:
\begin{equation}
\lbone_{\rm arm\,asym} \sim 6\ee{-19}/10^{-3} = 6\ee{-16}.
\end{equation}

\textbf{Channel 2 (gradient)}: For astrophysical sources at frequency
$f_\psi$, $\lambda_\psi = c/f_\psi$. At $f_\psi = 100$~Hz:
$\lambda_\psi = 3\ee{6}$~m, $L/\lambda_\psi = 1.3\ee{-3}$, so
$\sin^2(\pi L/\lambda_\psi) \sim (\pi L/\lambda_\psi)^2 \sim 1.7\ee{-5}$.
The bound:
\begin{equation}
\lbone_{\rm gradient} \sim 6\ee{-19}/1.7\ee{-5} = 3.5\ee{-14}.
\end{equation}

\textbf{Channel 3 (coating Process~A)}: $\delta\phi/\phi = (\lam-1)\cdot
\delta\psi\cdot d_{\rm coat}/L_{\rm arm} \sim (\lam-1)\cdot\delta\psi
\cdot 5\ee{-6}/(4\ee{3}) = (\lam-1)\cdot\delta\psi\cdot 1.25\ee{-9}$.
This is doubly suppressed by both $(\lam-1)$ and the coating/arm ratio.
Negligible.

\textbf{W3i7 corrected LIGO reach}:
\begin{equation}
\boxed{
\lbone^{\rm LIGO}_{\rm corrected} \sim 3.5\ee{-14}
\quad\text{(gradient channel, 100~Hz, O3 sensitivity)}
}
\label{eq:ligo_corrected}
\end{equation}

This is a \textbf{correction of 5 orders of magnitude} relative to
W3i5/W3i6. The $6\ee{-19}$ figure is \textbf{invalidated}.
[\textbf{CORRECTED}]
\end{finding}

%=============================================================================
\section{Complete Systematic Error Budget}
\label{sec:ligo_systematics}
%=============================================================================

\subsection{Classification of systematics}

We categorise all systematic effects that could mimic or mask a
psi-field signal in LIGO data:

\begin{longtable}{@{}p{3.5cm}p{2.5cm}p{2cm}p{5cm}@{}}
\caption{LIGO 6th-channel systematic error budget.}
\label{tab:ligo_systematics}\\
\toprule
\textbf{Systematic} & \textbf{Magnitude} & \textbf{Frequency dep.} &
  \textbf{Distinguishable from $\psi$?} \\
\midrule
\endfirsthead
\toprule
\textbf{Systematic} & \textbf{Magnitude} & \textbf{Frequency dep.} &
  \textbf{Distinguishable from $\psi$?} \\
\midrule
\endhead
\bottomrule
\endlastfoot
%
Newtonian noise (gravity gradient) &
  $h \sim 10^{-20}$ at 10~Hz &
  $\propto f^{-4}$ &
  \textbf{NO}. Scalar gravity perturbation is indistinguishable from
  scalar $\psi$-perturbation at the same frequency. Dominant systematic.
  \\
\midrule
Seismic/ground motion &
  $h \sim 10^{-19}$ at 10~Hz (after isolation) &
  $\propto f^{-2}$ above isolation corner &
  Partially. Seismic is primarily differential-mode; scalar $\psi$ is
  common-mode. Distinguishable via Wiener filtering with seismometer
  array.
  \\
\midrule
Thermal coating noise &
  $h \sim 2\ee{-24}\,\text{Hz}^{-1/2}$ at 100~Hz &
  $\propto f^{-1/2}$ &
  Yes. Thermal noise is mirror-local, uncorrelated between detectors.
  Psi-signal is coherent across all mirrors.
  \\
\midrule
Quantum shot noise &
  $h \sim 10^{-24}\,\text{Hz}^{-1/2}$ at 100~Hz &
  White (high-$f$ limit) &
  Yes. Shot noise is uncorrelated between detectors. Multi-detector
  cross-correlation rejects it as $1/\sqrt{T_{\rm int}}$.
  \\
\midrule
Laser frequency noise &
  $<10^{-7}$~Hz/$\sqrt{\rm Hz}$ (stabilised) &
  $\propto 1/f$ &
  Partially. Frequency noise is common-mode in both arms and thus mimics
  a scalar perturbation. Requires independent laser frequency reference
  (e.g., molecular spectroscopy).
  \\
\midrule
Scattered light &
  Variable; ``glitch-like'' &
  Broad, impulsive &
  Partially. Scatter fringe noise is non-stationary and can mimic
  broadband transients. Vetoed by auxiliary photodiode monitors.
  \\
\midrule
Magnetic coupling (Schumann resonances) &
  $h \sim 10^{-24}$~Hz$^{-1/2}$ at 14.1~Hz harmonics &
  Narrow lines at $\sim 7.8, 14.1, 20.3$~Hz &
  \textbf{CRITICAL}. Schumann resonances are global EM modes;
  they produce correlated signals across distant detectors. A scalar
  $\psi$-signal from solar EM activity would also be correlated.
  Schumann must be independently monitored and subtracted.
  \\
\midrule
Calibration errors &
  $\sim 5\%$ amplitude, $3^\circ$ phase &
  Frequency-dependent transfer function &
  Yes, but degrades parameter estimation. Not a false-positive risk.
  \\
\midrule
EM--GW source confusion &
  Depends on source &
  Source-dependent &
  The key issue. A compact binary merger produces both GW and EM
  emission. The GW (tensor) signal is subtracted, but residuals could
  alias into a scalar channel search. Requires precise GW subtraction
  at $<10^{-3}$ level.
  \\
\bottomrule
\end{longtable}

\subsection{The dominant systematic: Newtonian noise}

\begin{theorem}[Newtonian Noise as the LIGO 6th-Channel Floor]
\label{thm:newtonian_noise}
Newtonian (gravity-gradient) noise arises from density fluctuations
in the ground, atmosphere, and local infrastructure. These produce
time-varying gravitational accelerations on the test masses:
\begin{equation}
a_{\rm NN} = \GN \int \frac{\delta\rho(\mathbf{r}',t)\,
(\mathbf{r}-\mathbf{r}')}{|\mathbf{r}-\mathbf{r}'|^3}\,d^3r'.
\end{equation}

The resulting strain is:
\begin{equation}
h_{\rm NN}(f) \sim \frac{2\GN\,\rho_{\rm ground}\,c_{\rm seismic}^2}
{(2\pi f)^2\,L}\,S_{\rm seismic}^{1/2}(f),
\end{equation}
where $S_{\rm seismic}^{1/2}(f)$ is the seismic displacement noise ASD.

At 10~Hz for Advanced LIGO: $h_{\rm NN} \sim 10^{-20}$~Hz$^{-1/2}$.
This Newtonian noise is a \emph{scalar} gravitational perturbation
that is \textbf{physically identical in character} to a DFD
$\psi$-field perturbation. Both produce isotropic metric fluctuations
at the test mass locations.

\textbf{Below $\sim 20$~Hz, Newtonian noise is indistinguishable from
a psi-field signal in a single detector.} This sets an absolute
systematic floor:
\begin{equation}
\boxed{
h_{\rm floor}^{\rm NN}(f) \sim 10^{-20}\,\left(\frac{10~\text{Hz}}{f}
\right)^2~\text{Hz}^{-1/2}.
}
\end{equation}

Converting to $\lbone$ via the gradient coupling channel:
\begin{equation}
\lbone_{\rm floor}^{\rm NN} \sim \frac{h_{\rm floor}}{2\sin^2(\pi L/\lambda_\psi)}
\sim \frac{10^{-20}}{2\times 1.7\ee{-5}} \sim 3\ee{-16}
\quad\text{at 10~Hz.}
\end{equation}

At 100~Hz, Newtonian noise drops to $\sim 10^{-24}$~Hz$^{-1/2}$:
\begin{equation}
\lbone_{\rm floor}^{\rm NN}(100~\text{Hz}) \sim \frac{10^{-24}}{2\times 1.7\ee{-3}}
\sim 3\ee{-22}.
\end{equation}

However, at 100~Hz, the gradient coupling $\sin^2(\pi L/\lambda_\psi)
\sim 1.7\ee{-3}$ is larger, so the actual bound including both effects:
$\lbone \sim 3\ee{-22}/(\text{gradient coupling})$. We must be careful
with the accounting. The correct expression:
\begin{equation}
\delta\psi_{\rm min}(f) = \frac{h_{\rm noise}(f)}{2\sin^2(\pi L/\lambda_\psi(f))}.
\end{equation}

At 100~Hz: $\delta\psi_{\rm min} = 10^{-24}/(2\times 1.7\ee{-3})
= 2.9\ee{-22}$~Hz$^{-1/2}$.

With integration time $T_{\rm int} = 1$~yr $= 3.15\ee{7}$~s and
bandwidth $\Delta f = 1$~Hz:
\begin{equation}
\delta\psi_{\rm det} = \frac{\delta\psi_{\rm min}}{\sqrt{T_{\rm int}}}
= \frac{2.9\ee{-22}}{\sqrt{3.15\ee{7}}}
= 5.2\ee{-26}.
\end{equation}

This is the minimum detectable $\delta\psi$ at 100~Hz with 1~year
integration. The corresponding $\lbone$ depends on the \emph{source}
of the $\psi$-perturbation. For a cosmological stochastic background:
$\delta\psi \sim (\lam-1)\,\psiEM^{\rm cosmo} \sim (\lam-1)\times 9.6\ee{-6}$.
Then:
\begin{equation}
(\lam-1)_{\rm min} \sim \frac{5.2\ee{-26}}{9.6\ee{-6}} = 5.4\ee{-21}.
\end{equation}

For a modulated astrophysical source (e.g., pulsar with $\delta\psi_{\rm pulsar}$
at known frequency): the bound depends on the source amplitude.

\textbf{Summary}: The LIGO 6th-channel reach with proper systematic
accounting is:
\begin{equation}
\boxed{
\lbone^{\rm LIGO}_{\rm realistic} \sim 10^{-14}\text{--}10^{-21},
}
\end{equation}
depending critically on: (a) the frequency band, (b) the source model,
(c) whether Newtonian noise subtraction is achieved.
\end{theorem}

\subsection{The Schumann resonance problem}
\label{ssec:schumann}

\begin{finding}[Schumann Resonances as Correlated Background]
\label{find:schumann}
The Schumann resonances (global electromagnetic resonances of the
Earth-ionosphere cavity at $\sim 7.83$, 14.1, 20.3, 26.4, 32.6~Hz)
present a unique challenge. These produce:
\begin{enumerate}[nosep]
\item Correlated magnetic field fluctuations at LIGO Hanford and
  Livingston (and Virgo, KAGRA).
\item Coupling to the interferometer through magnetic actuation
  of the test masses (magnets on the penultimate mass) and through
  Barkhausen noise in magnetic shielding.
\item The correlated nature means that cross-correlation of two
  detectors does \textbf{not} reject Schumann-induced signals.
\end{enumerate}

A DFD psi-signal from solar EM activity would also appear at low
frequencies and would also be correlated across detectors (because the
Sun illuminates both sites equally). \textbf{This creates a direct
degeneracy between Schumann-coupled EM noise and a genuine solar
psi-signal.}

\textbf{Mitigation}: Independent Schumann monitoring with dedicated
magnetometers at each site, followed by Wiener-filter subtraction.
The LIGO collaboration has demonstrated Schumann subtraction at the
factor-of-$\sim 5$ level. Achieving subtraction to $<10^{-3}$ is
an open R\&D problem.
\end{finding}

%=============================================================================
\section{Null Tests for the LIGO 6th Channel}
\label{sec:null_tests}
%=============================================================================

\begin{theorem}[Suite of LIGO Psi-Signal Null Tests]
\label{thm:null_tests}
The following null tests can distinguish a true psi-signal from
systematics:

\textbf{Null Test 1: Polarisation separation.}
A true tensor GW has $h_+ \neq 0, h_\times \neq 0$ with the
characteristic $\cos 2\alpha$ antenna pattern. A scalar $\psi$-mode
has an isotropic antenna pattern (no $\alpha$-dependence). With three
or more non-co-aligned detectors (Hanford, Livingston, Virgo),
the polarisation content can be decomposed. A scalar-only signal
with no tensor component is inconsistent with GR sources but
consistent with DFD psi.

\textbf{Power}: This test requires $\SNR \gtrsim 8$ in each
detector for reliable polarisation decomposition.

\textbf{Null Test 2: EM--GW correlation.}
DFD predicts that $\delta\psi$ is sourced by EM stress-energy,
not by mass quadrupole moments. Therefore:
\begin{itemize}[nosep]
\item A psi-signal should correlate with EM events (magnetar flares,
  solar flares, gamma-ray bursts) but \emph{not} with the GR GW
  signal from the same event (which is sourced by mass motion).
\item For a binary neutron star merger: the GR GW signal peaks at
  coalescence; the psi-signal should peak with the post-merger EM
  emission (kilonova, jet).
\item For a magnetar giant flare: no GW signal expected in GR; DFD
  predicts a psi-signal from the EM energy release.
\end{itemize}

\textbf{Power}: Strongest null test. An excess signal correlated
with magnetar flares and \emph{not} with inspiral GW would be
strong evidence for psi. Requires joint EM/GW event analysis.

\textbf{Null Test 3: Frequency dependence.}
The DFD psi-field propagates at $c_s = c/n_\psi$. In the deep
Newtonian regime ($|\nabla\psi|/a_0 \gg 1$, i.e., the lab/near-Earth
environment), $n_\psi \approx 1$ and $c_s \approx c$. The
dispersion relation is:
\begin{equation}
\omega^2 = c_s^2 k^2 + m_\psi^2 c^4/\hbar^2,
\end{equation}
where $m_\psi$ is the effective psi-field mass. In DFD,
$m_\psi = 0$ (massless scalar), so $\omega = c_s k$. This means:
\begin{itemize}[nosep]
\item No dispersion (all frequencies travel at $c$).
\item The arrival time of a psi-pulse and a GW-pulse from the same
  source should be identical.
\end{itemize}
If a scalar signal shows frequency-dependent arrival time (dispersion),
it is not a DFD psi-signal.

\textbf{Power}: Moderate. Rules out massive scalar field alternatives
but not DFD-specific.

\textbf{Null Test 4: Stochastic background anisotropy.}
A cosmological stochastic $\psi$-background would have specific
anisotropy determined by the EM energy density distribution of the
universe (different from the mass distribution that sources the GR
stochastic GW background). The angular power spectrum of a psi
background should trace the EM luminosity-weighted matter distribution,
which differs from the mass-weighted distribution at the
$\sim 10\%$ level.

\textbf{Power}: Weak. Requires detection of the stochastic background
itself, which is a higher bar than single-source detection.

\textbf{Null Test 5: Day/night modulation.}
The Sun is the strongest nearby EM source. A DFD psi-perturbation
from the Sun has a diurnal modulation as the detector rotates
relative to the Sun. The diurnal modulation amplitude depends on
the antenna pattern for scalar modes:
\begin{equation}
F_\psi(\hat{n}) = \tfrac{1}{2}\sin^2\zeta,
\end{equation}
where $\zeta$ is the angle between $\hat{n}$ and the detector normal.
For a ground-based detector, the Sun traces a predictable path
in the antenna-pattern coordinates, producing a characteristic
diurnal signature.

\textbf{Power}: Strong for solar-sourced psi. The diurnal signature
is specific and predictable. However, it is degenerate with
temperature-driven diurnal systematics (thermal expansion of the
vacuum tubes, temperature-dependent electronics noise).
Mitigation: compare summer/winter amplitudes (thermal systematics
vary seasonally; solar psi-signal scales with $1/r_{\rm Sun}^2$,
which varies by $\sim 7\%$ over the year).
\end{theorem}

%=============================================================================
\section{Realistic LIGO 6th-Channel Reach Assessment}
\label{sec:ligo_realistic}
%=============================================================================

\begin{tcolorbox}[colback=red!5!white, colframe=red!60!black,
  title=\textbf{W3i7: LIGO 6th-Channel Honest Assessment}]

\textbf{The W3i5/W3i6 figure of $\lbone \sim 6\ee{-19}$ is
INVALIDATED.} It failed to account for:
\begin{enumerate}[nosep]
\item The common-mode rejection of scalar signals in a differential
  interferometer (suppression $\sim 10^{-3}$--$10^{-5}$).
\item Newtonian noise as an irreducible scalar systematic.
\item The Schumann resonance degeneracy.
\item Source-model dependence.
\end{enumerate}

\textbf{Revised reach estimates}:

\begin{center}
\renewcommand{\arraystretch}{1.25}
\begin{tabular}{@{}llll@{}}
\toprule
\textbf{Scenario} & \textbf{Frequency} & \textbf{$\lbone$ reach} &
  \textbf{Confidence} \\
\midrule
Magnetar flare correlation & $50$--$300$~Hz & $\sim 3\ee{-14}$ &
  Medium (depends on flare rate) \\
Solar stochastic (1~yr integration) & $50$--$500$~Hz & $\sim 5\ee{-21}$ &
  Low (Schumann degeneracy) \\
Pulsar correlation (Crab) & $59.5$~Hz & $\sim 10^{-17}$ &
  Low (Newtonian noise floor) \\
O3 archival, broadband & $20$--$2000$~Hz & $\sim 10^{-14}$ &
  Medium (gradient channel) \\
O5 projected & $10$--$5000$~Hz & $\sim 10^{-16}$ &
  Conditional on NN subtraction \\
\bottomrule
\end{tabular}
\end{center}

\textbf{Systematic floor}: Without Newtonian noise subtraction, the
floor is $\lbone \sim 3\ee{-14}$ at 100~Hz. With factor-10 NN
subtraction (projected for LIGO A+ / Cosmic Explorer): $\sim 3\ee{-15}$.
With factor-100 subtraction (optimistic): $\sim 3\ee{-16}$.

\textbf{Best-case realistic reach}: $\lbone \sim 10^{-14}$
from O3 archival data, improving to $\sim 10^{-16}$ with O5.

\textbf{Strategic conclusion}: The LIGO 6th channel remains the
second-most-sensitive passive approach (after ESS Lund optimistic
at $10^{-14}$, and potentially comparable). The $10^{10}\times$
gap over SQMS claimed in W3i6 is reduced to $\sim 10^{5}\times$.
Still extremely valuable as a parallel low-cost track, but no
longer the ``silver bullet'' previously advertised.

[\textbf{CORRECTED: 5 orders of magnitude downward from W3i6.}]
\end{tcolorbox}

%=============================================================================
\section{Updated Detection Ranking After LIGO Correction}
\label{sec:updated_ranking}
%=============================================================================

\begin{table}[H]
\centering
\caption{W3i7 corrected definitive DFD lab signal ranking.}
\label{tab:w3i7_ranking}
\renewcommand{\arraystretch}{1.25}
\begin{tabular}{@{}clllcc@{}}
\toprule
\textbf{Rank} & \textbf{Experiment} & \textbf{Observable} &
  \textbf{$\lbone$ reach} & \textbf{Status} & \textbf{TRL} \\
\midrule
\rowcolor{rowhi}
1 & ESS Lund (optimistic) (P3/P11) & Cross-corr.\ $\delta Q$
  & $1.14\ee{-14}$ & Projected 2027 & 6 \\
\rowcolor{rowhi}
2 & LIGO 6th channel (P11) --- \textbf{CORRECTED} & Gradient coupling
  $\delta\phi$ & $\sim 3\ee{-14}$ & Archival & 9 \\
\rowcolor{rowhi}
3 & ESS Lund (realistic) (P3/P11) & Cross-corr.\ $\delta Q$
  & $10^{-12}$--$10^{-13}$ & Projected 2027 & 6 \\
\rowcolor{rowhi}
4 & SQMS passive (P11) & $Q_0$ precision
  & $1.56\ee{-9}$ & \textbf{Achieved} & 9 \\
\rowcolor{rowmed}
5 & SQMS 4-cavity Phase~I & Incoherent $Q_0$ avg.
  & $7.8\ee{-10}$ & Projected 2027--28 & 8 \\
\rowcolor{rowmed}
6 & SQMS next-gen (P11/P15) & $Q_0$ at Nb$_3$Sn
  & $\sim 10^{-10}$ & Projected 2026--27 & 7 \\
\rowcolor{rowmed}
7 & P10+P11 (256-cavity array) & Coherent cross-corr.
  & $\sim 10^{-11}$ & Projected & 4 \\
\rowcolor{rowlo}
8 & P2 MEMS+P11 ($\Qpsi=10^{15}$) & SQUID mode amplitude
  & $\sim 10^{-11}$ & Speculative & 2 \\
\rowcolor{rowlo}
9 & P2 MEMS+P11 ($\Qpsi=10^{9}$) & SQUID mode amplitude
  & $\sim 10^{-5}$ & Speculative & 2 \\
\rowcolor{rowlo}
10 & P1 timing (best lab source) & Non-reciprocal $\delta t$
  & $\sim 10^{4}$ & \textbf{Dead} & --- \\
\bottomrule
\end{tabular}
\end{table}

\begin{finding}[LIGO Drops from \#1 to \#2; ESS Lund Now Leads]
\label{find:ranking_change}
The LIGO 6th-channel correction moves it from the dominant position
($10^{-19}$) to approximate parity with ESS Lund optimistic
($\sim 10^{-14}$). ESS Lund now leads the ranking because:
\begin{enumerate}[nosep]
\item Its systematic environment is better understood (lock-in to
  beam pulsing at 14~Hz rejects uncorrelated noise).
\item It does not suffer from the common-mode rejection problem
  (cavity $\delta Q/Q$ is a single-detector measurement).
\item No Newtonian noise floor (SRF cavities are not gravitational
  wave detectors).
\end{enumerate}

The gap between passive and active remains enormous:
$\sim 10^{-14}$ vs $\sim 10^{-5}$ = 9 orders of magnitude (reduced
from 14 in W3i6 but still fundamental).
[\textbf{CORRECTED}]
\end{finding}


%=============================================================================
%=============================================================================
\part{SQMS Phase~I: Detailed Experimental Protocol}
\label{part:SQMS}
%=============================================================================
%=============================================================================

%=============================================================================
\section{Phase~I Configuration: 4-Cavity Incoherent Array}
\label{sec:phase1_config}
%=============================================================================

\subsection{Hardware specification}

\begin{table}[H]
\centering
\caption{SQMS Phase~I hardware specification.}
\label{tab:phase1_hardware}
\renewcommand{\arraystretch}{1.2}
\begin{tabular}{@{}lll@{}}
\toprule
\textbf{Component} & \textbf{Specification} & \textbf{Source} \\
\midrule
Cavities & 4$\times$ TESLA 9-cell, 1.3~GHz, N-doped Nb &
  Fermilab SQMS inventory \\
Cryostat & Existing dilution refrigerator (DR) or 2~K He bath &
  Fermilab SQMS \\
$T_{\rm op}$ & Primary: 20~mK (DR); Secondary: 1.8~K (He bath) &
  DR for best $Q_0$ \\
$Q_0$ target & $\geq 2\ee{10}$ (1.8~K); $\geq 2\ee{11}$ (20~mK) &
  Demonstrated at SQMS \\
RF source & Continuous-wave (CW) at 1.3~GHz, phase-locked &
  Existing LLRF \\
$E_{\rm acc}$ & Low-gradient: $E_{\rm acc} \leq 5$~MV/m &
  Below medium-field $Q$-slope \\
Readout & 4$\times$ independent SQUID or RF pickup readout &
  Upgrade needed \\
Shielding & $\mu$-metal: $B_{\rm residual} < 1$~nT per cavity &
  Standard SQMS practice \\
Vibration isolation & Passive + active ($<10^{-9}$~m/$\sqrt{\rm Hz}$
  at 1~Hz) & Standard cryostat isolation \\
Temperature stability & $\delta T/T < 10^{-4}$ at 20~mK over $10^6$~s &
  DR stability spec \\
\bottomrule
\end{tabular}
\end{table}

\subsection{Cavity placement and systematics separation}

\begin{finding}[Cavity Arrangement for Systematic Diagnosis]
\label{find:cavity_arrangement}
The four cavities should be arranged in a $2\times 2$ configuration:

\begin{enumerate}[nosep]
\item \textbf{Pair A} (cavities 1,2): Same cryostat, same vacuum
  space, $\sim 0.5$~m separation. Shares all environmental
  systematics (temperature, magnetic field, vibration).
\item \textbf{Pair B} (cavities 3,4): Different cryostat (or same
  cryostat but different vacuum space), $>2$~m separation. Independent
  environmental systematics.
\end{enumerate}

This arrangement enables three cross-correlations:
\begin{itemize}[nosep]
\item \textbf{Intra-pair A} ($C_{12}$): tests for correlated
  environmental effects.
\item \textbf{Intra-pair B} ($C_{34}$): same, independent sample.
\item \textbf{Inter-pair} ($C_{13}, C_{14}, C_{23}, C_{24}$): tests
  for genuine common-mode signals (psi-field or other global effects).
\end{itemize}

A genuine psi-field signal produces $C_{12} = C_{34} = C_{13} = C_{14}$
(all correlations equal). Environmental systematics produce
$C_{12} \gg C_{13}$ (intra-pair correlation dominates).
\end{finding}

%=============================================================================
\section{Measurement Protocol}
\label{sec:measurement_protocol}
%=============================================================================

\subsection{Phase~I-A: Baseline characterisation (Months 1--6)}

\begin{enumerate}[label=\textbf{Step \arabic*.}, itemsep=6pt]

\item \textbf{Individual cavity $Q_0$ mapping.}
  For each cavity independently:
  \begin{itemize}[nosep]
  \item Measure $Q_0(E_{\rm acc})$ at 5 gradient points:
    $E_{\rm acc} = 1, 2, 3, 4, 5$~MV/m.
  \item Measure $Q_0(T)$ at 8 temperature points:
    $T = 20, 50, 100, 200, 500$~mK and $1.2, 1.5, 1.8$~K.
  \item Measure $Q_0(B_{\rm ext})$ at 5 magnetic field values:
    $B_{\rm ext} = 0, 0.5, 1, 2, 5$~nT (varied with external coils).
  \item Duration: $\sim 2$~weeks per cavity, $\sim 2$~months total.
  \end{itemize}

  \textbf{Purpose}: Establish the BCS surface resistance model
  $R_s^{\rm BCS}(T, B, E)$ for each cavity. Any deviation from the
  standard BCS + residual resistance model at the lowest temperatures
  is a potential psi-channel signature.

\item \textbf{Multi-frequency $Q_0$ measurements.}
  Each TESLA 9-cell cavity supports multiple TM and TE modes
  at different frequencies. Measure $Q_0$ at:
  \begin{itemize}[nosep]
  \item $\text{TM}_{010}$: 1.3~GHz (fundamental accelerating mode)
  \item $\text{TM}_{011}$: $\sim 2.4$~GHz
  \item $\text{TE}_{011}$: $\sim 1.7$~GHz
  \item Higher-order modes as accessible
  \end{itemize}

  \textbf{Purpose}: The psi-channel loss has a specific frequency
  dependence (Eq.~(\ref{eq:psi_loss_freq}) below), while BCS loss
  scales as $\omega^2$ and residual loss is frequency-independent.
  Multi-frequency measurements separate these contributions.

\item \textbf{Environmental monitoring baseline.}
  Record 30~days of continuous environmental data:
  \begin{itemize}[nosep]
  \item Temperature: $\delta T(t)$ at each cavity with $<1~\mu$K resolution.
  \item Magnetic field: 3-axis fluxgate at each cavity, $< 10$~pT resolution.
  \item Vibration: 3-axis accelerometer on each cryostat, $<10^{-10}$~g/$\sqrt{\rm Hz}$.
  \item Pressure: He bath pressure, $< 1~\mu$bar resolution.
  \end{itemize}

  \textbf{Purpose}: Establish the transfer functions
  $\partial Q_0/\partial T$, $\partial Q_0/\partial B$,
  $\partial Q_0/\partial x$ for each cavity. These are used
  to subtract known environmental contributions from the $Q_0(t)$
  time series.
\end{enumerate}

\subsection{Phase~I-B: Science data taking (Months 7--18)}

\begin{enumerate}[label=\textbf{Step \arabic*.}, resume, itemsep=6pt]

\item \textbf{Continuous $Q_0(t)$ monitoring.}
  Operate all 4 cavities simultaneously at fixed operating point
  ($E_{\rm acc} = 2$~MV/m, $T = 20$~mK, $B < 1$~nT) for
  $\geq 10^6$~s ($\sim 12$~days) per science run. Target: 10 science
  runs over 12 months.

  Readout: cavity bandwidth $\Delta f_{\rm cav} = f_0/Q_0
  = 1.3\ee{9}/(2\ee{11}) = 6.5$~mHz. Ring-down time
  $\tau = Q_0/(\pi f_0) = 2\ee{11}/(4.1\ee{9}) = 49$~s.
  Sample $Q_0$ by measuring the ring-down time constant every
  $\sim 5\tau = 245$~s.

  \textbf{Per-sample precision}: The fractional $Q_0$ precision from
  a single ring-down is:
  \begin{equation}
  \frac{\sigma_Q}{Q_0}\bigg|_{\rm single} \sim \frac{1}{\sqrt{2\pi f_0\,\tau}}
  \cdot\frac{1}{\SNR_{\rm readout}} \sim \frac{1}{\sqrt{2\pi\times 1.3\ee{9}
  \times 49}}\cdot\frac{1}{10^3} \sim \frac{1}{6.3\ee{5}\times 10^3}
  \sim 1.6\ee{-9}.
  \end{equation}

  Over $N_{\rm samples} = 10^6/245 \approx 4000$ samples per run:
  \begin{equation}
  \frac{\sigma_Q}{Q_0}\bigg|_{\rm run} = \frac{1.6\ee{-9}}{\sqrt{4000}}
  = 2.5\ee{-11}.
  \end{equation}

  This is well below the current SQMS bound ($\delta Q/Q \sim 10^{-3}$),
  suggesting the current bound is dominated by systematics, not
  statistical precision. The Phase~I programme must identify and
  subtract these systematics.

\item \textbf{Cross-correlation analysis.}
  For each pair of cavities $(i,j)$, compute:
  \begin{equation}
  C_{ij}(\tau) = \frac{1}{T}\int_0^T \frac{\delta Q_i(t)}{Q_{0,i}}
  \cdot\frac{\delta Q_j(t+\tau)}{Q_{0,j}}\,dt,
  \end{equation}
  where $\delta Q_i(t) = Q_i(t) - \langle Q_i\rangle$ and $\tau$ is
  the time lag.

  \textbf{Expected signatures}:
  \begin{itemize}[nosep]
  \item \textbf{Psi-field}: $C_{ij}(0) > 0$ for all pairs, with
    amplitude $\propto (\lam-1)^2$. Time lag $\tau_\psi = d_{ij}/c_s
    \sim 0.5~\text{m}/(3\ee{8}~\text{m/s}) \sim 2$~ns (unresolvable
    at 245~s sampling; appears as $\tau = 0$).
  \item \textbf{Temperature drift}: $C_{12}(0) > 0, C_{34}(0) > 0$
    (intra-pair correlated), but $C_{13}(0) \approx 0$ if cryostats
    are thermally independent.
  \item \textbf{Magnetic field}: Similar to temperature, but may be
    partially correlated across cryostats (Schumann resonances,
    building-scale fields).
  \item \textbf{Cosmic rays}: Produces correlated quasiparticle bursts
    in closely spaced cavities ($C_{12}$) but uncorrelated at large
    separations ($C_{13}$).
  \end{itemize}

\item \textbf{Environmental subtraction.}
  Construct the environmentally cleaned $Q_0$ time series:
  \begin{equation}
  \tilde{Q}_i(t) = Q_i(t) - \sum_k \alpha_{ik}\,E_k(t),
  \end{equation}
  where $E_k(t)$ are the environmental monitor channels (temperature,
  magnetic field, vibration, pressure) and $\alpha_{ik}$ are the
  transfer coefficients determined in Phase~I-A. The cleaned
  cross-correlation $\tilde{C}_{ij}$ should be consistent with zero
  if no psi-signal is present.
\end{enumerate}

%=============================================================================
\section{Detection vs Null Criteria}
\label{sec:detection_criteria}
%=============================================================================

\begin{theorem}[Phase~I Detection Criterion]
\label{thm:detection_criterion}
A DFD psi-field detection is claimed if and only if:
\begin{enumerate}[nosep]
\item The environmentally cleaned inter-pair cross-correlation is
  non-zero:
  $\tilde{C}_{13}(0) > 3\sigma_{\tilde{C}}$ (discovery threshold at
  $3\sigma$; $5\sigma$ for ``observation'').
\item The intra-pair and inter-pair correlations are \emph{consistent}:
  $\tilde{C}_{12}(0) \approx \tilde{C}_{13}(0)$ within statistical
  uncertainties (i.e., the signal is not dominated by local
  environmental effects).
\item The frequency dependence of $\delta Q/Q$ across modes (TM$_{010}$,
  TM$_{011}$, TE$_{011}$) is consistent with the DFD prediction
  (Eq.~(\ref{eq:psi_loss_freq})) and \emph{inconsistent} with BCS or
  residual resistance scaling.
\item At least one independent systematic check passes: the correlation
  is absent when one cavity is detuned by $>10$ bandwidths (psi
  coupling is resonant; environmental coupling is broadband).
\end{enumerate}

If all four criteria are met simultaneously, the result constitutes
evidence for DFD psi-coupling at the measured $\lbone$ value.
\end{theorem}

\begin{finding}[Null Result Interpretation]
\label{find:null_result}
A null result ($\tilde{C}_{13}(0)$ consistent with zero at
$3\sigma$) after the full Phase~I programme establishes:
\begin{equation}
\lbone < \frac{1.56\ee{-9}}{\sqrt{N_{\rm eff}}},
\end{equation}
where $N_{\rm eff}$ is the effective number of independent
measurements after systematic subtraction. For 4 cavities with
10 science runs:
\begin{itemize}[nosep]
\item Best case ($N_{\rm eff} = 4\times 10 = 40$): $\lbone < 2.5\ee{-10}$.
\item Realistic (systematics reduce $N_{\rm eff}$ to $\sim 4$):
  $\lbone < 7.8\ee{-10}$.
\item Conservative ($N_{\rm eff} = 2$, strong systematics):
  $\lbone < 1.1\ee{-9}$.
\end{itemize}

\textbf{Even the conservative null result improves the SQMS bound}
by $\sim 40\%$, confirming the methodology.

A null result also provides:
\begin{enumerate}[nosep]
\item The measured systematic floor for SRF cavity $Q_0$
  cross-correlations (input to Phase~II/III planning).
\item Transfer coefficients for environmental coupling (reusable
  across all future SQMS measurements).
\item Bounds on correlated non-DFD anomalies (cosmic ray correlations,
  unknown material effects).
\end{enumerate}
\end{finding}

%=============================================================================
\section{The Psi-Channel Loss Frequency Dependence}
\label{sec:freq_dependence}
%=============================================================================

The DFD-predicted psi-channel contribution to cavity loss is:
\begin{equation}
\frac{1}{Q_\psi(\omega)} = \frac{(\lam-1)^2\,\Heff^2(\omega)}{\Qpsi}
\cdot\frac{4\pi\GN\,u_{\rm EM}}{\muG\,c^2}
\cdot\frac{\Opsi(\omega)}{\omega},
\label{eq:psi_loss_freq}
\end{equation}
where $\Opsi(\omega)$ is the psi-mode frequency matched to the EM mode
at $\omega$. In the DFD framework, the psi-mode frequency tracks the EM
mode: $\Opsi = \omega$ (resonant coupling). Therefore:
\begin{equation}
\frac{1}{Q_\psi(\omega)} = \frac{(\lam-1)^2\,\Heff^2(\omega)}{\Qpsi}
\cdot\frac{4\pi\GN\,u_{\rm EM}}{\muG\,c^2}.
\end{equation}

The frequency dependence enters through $\Heff^2(\omega)$, the
overlap integral between the EM mode and the psi-mode. For a TESLA
cavity:
\begin{itemize}[nosep]
\item TM$_{010}$ (1.3~GHz): $\Heff = 3\ee{-5}$ (baseline).
\item TM$_{011}$ (2.4~GHz): $\Heff \approx 2.1\ee{-5}$ (different
  field distribution; more concentrated on axis).
\item TE$_{011}$ (1.7~GHz): $\Heff \approx 1.5\ee{-5}$ (TE mode has
  maximum $E$ at the equator, different overlap).
\end{itemize}

The ratio of psi-channel losses at two frequencies:
\begin{equation}
\frac{1/Q_\psi(\omega_1)}{1/Q_\psi(\omega_2)} = \frac{\Heff^2(\omega_1)}
{\Heff^2(\omega_2)}.
\end{equation}

This is \emph{distinct} from BCS scaling ($R_s^{\rm BCS} \propto \omega^2$)
and residual resistance (frequency-independent). The multi-frequency
measurement directly tests the DFD prediction.

\begin{finding}[Multi-Frequency Diagnostic Power]
\label{find:multi_freq}
Measuring $Q_0$ at three frequencies provides two independent ratios
that can be compared against three models:

\begin{center}
\renewcommand{\arraystretch}{1.2}
\begin{tabular}{@{}lll@{}}
\toprule
\textbf{Loss model} & $Q_0^{-1}(2.4~\text{GHz})/Q_0^{-1}(1.3~\text{GHz})$ &
  $Q_0^{-1}(1.7~\text{GHz})/Q_0^{-1}(1.3~\text{GHz})$ \\
\midrule
BCS ($\propto \omega^2$) & $(2.4/1.3)^2 = 3.41$ & $(1.7/1.3)^2 = 1.71$ \\
Residual (const.) & 1.0 & 1.0 \\
Psi ($\propto \Heff^2$) & $(2.1/3.0)^2 = 0.49$ & $(1.5/3.0)^2 = 0.25$ \\
\bottomrule
\end{tabular}
\end{center}

The psi-channel signature is dramatically different from BCS
($0.49$ vs $3.41$) and from residual ($0.49$ vs $1.0$). A
measurement of the multi-frequency $Q_0$ ratio to $\sim 10\%$
precision would unambiguously identify or exclude a psi-channel
contribution.

\textbf{This is the most powerful single diagnostic in the Phase~I
programme.} [\textbf{NEW}]
\end{finding}


%=============================================================================
%=============================================================================
\part{Indirect $\Qpsi$ Constraints from Existing Data}
\label{part:Qpsi}
%=============================================================================
%=============================================================================

%=============================================================================
\section{Can Existing SQMS Data Already Constrain $\Qpsi$?}
\label{sec:existing_data}
%=============================================================================

\subsection{The fundamental constraint equation}

From Eq.~(\ref{eq:psi_loss_freq}), the psi-channel contribution to the
inverse quality factor is:
\begin{equation}
\frac{1}{Q_\psi} = \frac{(\lam-1)^2}{\Qpsi}\,\Heff^2
\,\frac{4\pi\GN\,u_{\rm EM}}{\muG\,c^2}.
\label{eq:Qpsi_constraint_basic}
\end{equation}

The measured residual loss (after BCS subtraction) provides an upper
bound:
\begin{equation}
\frac{1}{Q_\psi} < \frac{1}{Q_{\rm res}^{\rm meas}},
\end{equation}
where $Q_{\rm res}^{\rm meas}$ is the measured residual quality factor
(the $Q_0$ at $T \to 0$ after BCS subtraction).

This gives the joint constraint:
\begin{equation}
\boxed{
\frac{(\lam-1)^2}{\Qpsi} < \frac{1}{Q_{\rm res}\,\Heff^2}
\cdot\frac{\muG\,c^2}{4\pi\GN\,u_{\rm EM}}.
}
\label{eq:joint_constraint}
\end{equation}

\subsection{Numerical evaluation with SQMS data}

The SQMS programme has published extensive $Q_0(T)$ data for N-doped
Nb cavities. The key numbers:

\begin{table}[H]
\centering
\caption{SQMS published parameters for $\Qpsi$ constraint.}
\label{tab:sqms_params}
\renewcommand{\arraystretch}{1.2}
\begin{tabular}{@{}lll@{}}
\toprule
\textbf{Parameter} & \textbf{Value} & \textbf{Source} \\
\midrule
$Q_{\rm res}$ (residual, $T \to 0$) & $\sim 5\ee{10}$--$2\ee{11}$
  & Fermilab SQMS 2023--2025 \\
$\Heff$ & $3\ee{-5}$ & TESLA geometry calculation \\
$u_{\rm EM}$ at $E_{\rm acc} = 5$~MV/m & $\sim 10^4$~J/m$^3$ &
  Standard SRF \\
$\muG$ & 0.657 & SPARC galaxy fits (W3i5) \\
\bottomrule
\end{tabular}
\end{table}

Inserting into Eq.~(\ref{eq:joint_constraint}):
\begin{align}
\frac{(\lam-1)^2}{\Qpsi}
&< \frac{1}{5\ee{10}\times (3\ee{-5})^2}
\cdot\frac{0.657\times 9\ee{16}}{4\pi\times 6.674\ee{-11}\times 10^4}
\notag\\
&= \frac{1}{5\ee{10}\times 9\ee{-10}}
\cdot\frac{5.91\ee{16}}{8.39\ee{-6}}
\notag\\
&= \frac{1}{4.5\ee{1}}
\cdot 7.04\ee{21}
\notag\\
&= 1.56\ee{20}.
\label{eq:joint_numerical}
\end{align}

Using the SQMS bound $\lbone < 1.56\ee{-9}$, hence
$(\lam-1)^2 < 2.43\ee{-18}$:
\begin{equation}
\Qpsi > \frac{2.43\ee{-18}}{1.56\ee{20}} \times Q_{\rm res}
\sim \frac{2.43\ee{-18}}{1.56\ee{20}} \approx 1.6\ee{-38}.
\end{equation}

\begin{finding}[$\Qpsi$ Constraint from Existing SQMS Data: Trivial]
\label{find:Qpsi_trivial}
The existing SQMS data provide only the trivially weak constraint
$\Qpsi > 1.6\ee{-38}$, which is physically meaningless (quality
factors are positive and typically $\gg 1$).

The reason: the constraint equation involves the \emph{product}
$(\lam-1)^2/\Qpsi$, and with only one observable ($Q_{\rm res}$),
the two unknowns $\lbone$ and $\Qpsi$ are completely degenerate.
The SQMS bound on $\lbone$ was derived \emph{assuming}
$\Qpsi = 10^9$. Using the resulting $\lbone$ bound to then
constrain $\Qpsi$ is circular.

\textbf{To break the degeneracy, one needs at least two independent
observables with different $(\lam-1)^2/\Qpsi$ scalings.}
[\textbf{VERIFIED} --- inherits from W3i6 Finding 7.4]
\end{finding}

%=============================================================================
\section{Breaking the Degeneracy: Multi-Observable Strategy}
\label{sec:breaking_degeneracy}
%=============================================================================

\subsection{Strategy 1: Multi-frequency residual $Q_0$}

If the residual loss at two different EM frequencies has different
$\Heff$ values, the ratio of residual losses constrains the psi-channel
contribution independently of both $\lbone$ and $\Qpsi$:

\begin{equation}
\frac{1/Q_{\rm res}(\omega_1) - 1/Q_{\rm res}^{\rm mat}(\omega_1)}
{1/Q_{\rm res}(\omega_2) - 1/Q_{\rm res}^{\rm mat}(\omega_2)}
= \frac{\Heff^2(\omega_1)}{\Heff^2(\omega_2)},
\label{eq:multi_freq_ratio}
\end{equation}
where $Q_{\rm res}^{\rm mat}$ is the material residual (from trapped
flux, surface impurities, etc., assumed to be frequency-independent
or with known frequency dependence).

If this ratio is measured and found to be consistent with the DFD
prediction ($\Heff^2(\omega_1)/\Heff^2(\omega_2)$), then the absolute
value of either residual excess gives $(\lam-1)^2/\Qpsi$ directly:
\begin{equation}
\frac{(\lam-1)^2}{\Qpsi} = \left[\frac{1}{Q_{\rm res}(\omega_1)}
- \frac{1}{Q_{\rm res}^{\rm mat}(\omega_1)}\right]
\cdot\frac{1}{\Heff^2(\omega_1)}\cdot\frac{\muG\,c^2}
{4\pi\GN\,u_{\rm EM}}.
\label{eq:Qpsi_multi_freq}
\end{equation}

\textbf{Obstacle}: The material residual $Q_{\rm res}^{\rm mat}$ is
itself uncertain at the $\sim 10^{-11}$ level. Distinguishing a
psi-channel contribution of $1/Q_\psi \sim 10^{-11}$ from the
material residual requires:
\begin{equation}
\sigma(1/Q_{\rm res}^{\rm mat}) < 10^{-11}.
\end{equation}
Current material science understanding of trapped-flux residual
resistance gives $\sigma \sim 10^{-11}$--$10^{-10}$. The psi-channel
contribution is at the noise floor of material characterisation.

\subsection{Strategy 2: Temperature-dependent residual $Q_0$}

\begin{finding}[Temperature Signature of Psi-Channel Loss]
\label{find:temperature_signature}
The BCS surface resistance has well-characterised temperature
dependence:
\begin{equation}
R_s^{\rm BCS}(T) = A\,\omega^2\,\frac{e^{-\Delta/(k_{\rm B}T)}}{T},
\end{equation}
while the psi-channel loss is temperature-independent (it depends on
$\Meff$, $\Qpsi$, $\lbone$, and $\Heff$, none of which depend on $T$).

At $T \ll T_c/10 \approx 0.9$~K for Nb, the BCS contribution is
exponentially suppressed. The measured $Q_0^{-1}(T)$ curve should
flatten to a constant (the sum of material residual and psi-channel):
\begin{equation}
Q_0^{-1}(T \to 0) = Q_{\rm res,mat}^{-1} + Q_\psi^{-1}.
\end{equation}

If this ``residual plateau'' has a component that varies with
$u_{\rm EM}$ (field gradient), it could be the psi-channel:
$Q_\psi^{-1} \propto u_{\rm EM}$, while $Q_{\rm res,mat}^{-1}$ is
field-independent below the medium-field $Q$-slope onset.

\textbf{The decisive measurement}: Compare the residual $Q_0$ at
$E_{\rm acc} = 1$~MV/m and $E_{\rm acc} = 5$~MV/m in the same cavity
at $T = 20$~mK:
\begin{equation}
\Delta Q_0^{-1} = Q_0^{-1}(5~\text{MV/m}) - Q_0^{-1}(1~\text{MV/m}).
\end{equation}

If $\Delta Q_0^{-1} > 0$ and scales as $E^2$ (i.e., as $u_{\rm EM}$),
it is consistent with a psi-channel. If $\Delta Q_0^{-1} = 0$, the
psi-channel is below the measurement floor.

\textbf{Current SQMS data}: Existing $Q_0$ vs $E_{\rm acc}$ curves at
low temperatures show a ``low-field $Q$-slope'' (increasing $Q_0$ with
$E_{\rm acc}$ up to $\sim 5$~MV/m) in N-doped cavities. This is
generally attributed to TLS saturation: at higher fields, TLS defects
are driven out of resonance, reducing their loss contribution. The
psi-channel would produce the \emph{opposite} trend (decreasing $Q_0$
with increasing $E_{\rm acc}$).

\textbf{Existing data already weakly constrain the psi-channel}:
since $Q_0$ \emph{increases} with $E_{\rm acc}$ in the low-field
regime, any psi-channel must produce a loss smaller than the
TLS saturation gain. This implies:
\begin{equation}
Q_\psi^{-1}(5~\text{MV/m}) < Q_{\rm TLS}^{-1}(1~\text{MV/m})
- Q_{\rm TLS}^{-1}(5~\text{MV/m}) \sim 10^{-11}.
\end{equation}

From this: $(\lam-1)^2/\Qpsi < 10^{-11}/(A_{\rm geo}) \sim 10^{9}$
(using $A_{\rm geo} = \Heff^2 \cdot 4\pi G u_{\rm EM}/(\muG c^2)
\sim 10^{-20}$). This gives $(\lam-1)^2/\Qpsi < 10^{9} \times 10^{-20}
= 10^{-11}$, hence:
\begin{equation}
\frac{(\lam-1)^2}{\Qpsi} < \frac{10^{-11}}{1} = 10^{-11}.
\end{equation}

With $\lbone < 1.56\ee{-9}$:
\begin{equation}
\Qpsi > \frac{(1.56\ee{-9})^2}{10^{-11}} = \frac{2.43\ee{-18}}{10^{-11}}
= 2.4\ee{-7}.
\end{equation}

Still trivially weak. The problem is that $10^{-11}$ is a very loose
bound on $Q_\psi^{-1}$.
\end{finding}

\subsection{Strategy 3: Dedicated dual-cavity differential measurement}

\begin{theorem}[Breaking the $\lbone/\Qpsi$ Degeneracy]
\label{thm:breaking_degeneracy}
The only way to separately constrain $\lbone$ and $\Qpsi$ is through
two observables with \emph{different} scaling:

\textbf{Observable 1}: Passive $Q_0$ bound scales as
$\delta Q/Q \propto (\lam-1)^2/\Qpsi$.

\textbf{Observable 2}: Psi-mode ringdown time scales as
$\tau_\psi = \Qpsi/\Opsi$.

If the psi-mode can be excited (by an intense SRF source) and its
ringdown observed (in a separate detector cavity), $\Qpsi$ is measured
directly. But exciting the psi-mode requires closing the active
detection gap, which requires knowing $\Qpsi$ --- a circular problem.

\textbf{The degeneracy cannot be broken with passive measurements
alone.} Only if a psi-signal is detected (in any channel) can $\Qpsi$
be extracted from the signal properties (amplitude, coherence time,
spectral shape).

\textbf{Practical path}: If Phase~I or ESS Lund detects a non-zero
$\tilde{C}_{ij}(0)$, the correlation timescale
$\tau_{\rm corr} = \int C(\tau)\,d\tau / C(0)$ gives
$\tau_\psi \sim \Qpsi/\Opsi$. For $\Opsi = 8.2\ee{9}$~rad/s:
\begin{equation}
\Qpsi \approx \Opsi\,\tau_{\rm corr}.
\end{equation}
A correlation timescale of $1~\mu$s gives $\Qpsi \sim 8200$;
of 1~s gives $\Qpsi \sim 8.2\ee{9}$; of 1~hr gives
$\Qpsi \sim 3\ee{13}$.
\end{theorem}

%=============================================================================
\section{Existing Datasets for Reanalysis}
\label{sec:existing_datasets}
%=============================================================================

\begin{finding}[Datasets That Should Be Reanalysed for Indirect $\Qpsi$ Constraints]
\label{find:datasets}

\begin{enumerate}[label=\textbf{D\arabic*.}, itemsep=6pt]

\item \textbf{Fermilab SQMS: Multi-frequency $Q_0$ data (2019--2025).}
  The SQMS programme has measured $Q_0$ at multiple frequencies in
  TESLA, elliptical, and spoke cavities. The multi-frequency residual
  ratio (Eq.~\ref{eq:multi_freq_ratio}) can be computed from
  published data. If the ratio is inconsistent with the DFD prediction,
  it provides a null result. If consistent, it constrains
  $(\lam-1)^2/\Qpsi$.

  \textbf{Key dataset}: Romanenko et al., ``Ultra-high quality factors
  in superconducting niobium cavities in ambient magnetic fields,''
  Phys.\ Rev.\ Applied (2020). Contains $Q_0(f)$ at 1.3 and 2.6~GHz
  in the same cavity geometry.

\item \textbf{DESY: TESLA cavity $Q_0(T)$ database (2000--2020).}
  Two decades of $Q_0(T)$ curves for hundreds of TESLA cavities.
  The residual plateau statistics across many cavities constrain the
  cavity-to-cavity variance of $Q_{\rm res}$. If the psi-channel
  contributes uniformly to all cavities, it shows up as a universal
  residual floor that is independent of surface preparation.

  \textbf{Key question}: Is there a ``universal residual'' component
  in the DESY database that cannot be explained by trapped flux or
  surface defects?

\item \textbf{JLab: C100 cryomodule $Q_0$ correlations (2015--2025).}
  The CEBAF C100 cryomodules contain 8 cavities per module, operating
  simultaneously. If correlated $Q_0$ fluctuations are present in
  JLab operational data, they could be a psi-field signature. The
  existing LLRF monitoring data is extensive but has not been
  analysed for inter-cavity correlations at the precision needed.

  \textbf{Key question}: Do C100 cavities within a module show
  correlated $Q_0$ fluctuations above the environmental background?

\item \textbf{LIGO: Cavity length stability data (2015--2025).}
  The LIGO Fabry--Perot arm cavities are exquisite frequency
  references. The long-term frequency drift data ($\delta f/f$ over
  days to months) may contain a component from psi-field fluctuations.
  This is a different observable from the strain channel: it measures
  the cavity frequency directly, not the differential arm length.

  \textbf{Caveat}: LIGO cavities are not superconducting; the
  psi-coupling via Process~A is suppressed by the much lower EM
  energy density in optical cavities compared with SRF cavities.

\item \textbf{PTB Braunschweig: Sapphire oscillator frequency stability
  (2010--2025).}
  The cryogenic sapphire oscillators at PTB achieve fractional
  frequency stability $\sigma_y \sim 10^{-16}$ at 1~s averaging time.
  Any psi-field coupling would appear as an anomalous instability
  component. The sapphire oscillators operate at $\sim 10$~GHz with
  $Q_0 \sim 10^9$ at 4~K. The psi-coupling is weaker (lower $Q_0$
  than SRF) but the frequency stability measurement is much more
  precise.
\end{enumerate}
\end{finding}


%=============================================================================
%=============================================================================
\part{Patent 1: Definitive Final Assessment}
\label{part:P1_final}
%=============================================================================
%=============================================================================

%=============================================================================
\section{P1 Correction History}
\label{sec:P1_history}
%=============================================================================

\begin{table}[H]
\centering
\caption{Complete correction history for Patent 1 (Field Drag /
Gravitational Cloaking) across all iterations.}
\label{tab:P1_history}
\renewcommand{\arraystretch}{1.25}
\begin{tabular}{@{}p{1.8cm}p{4cm}p{4cm}p{3.5cm}@{}}
\toprule
\textbf{Iteration} & \textbf{Key finding} & \textbf{Impact on P1} &
  \textbf{Status after} \\
\midrule
W3i1 & Initial analysis of drag timing & 4.0~ns signal claimed &
  \textcolor{strongcolor}{Promising} \\
W3i2 & Cross-referenced with P11, P2 & 4.0~ns confirmed;
  linked to SRF source & \textcolor{strongcolor}{Promising} \\
W3i3 & Cosmological corrections begun & Minor adjustments &
  \textcolor{strongcolor}{Stable} \\
W3i4 & Full numerical calculations & $u_{\rm EM}$ requirements
  quantified & \textcolor{conjcolor}{Caution} \\
W3i5 & EM-only sourcing: $\psicosmo^{\rm EM} = 1.3\ee{-7}$;
  $\Meff$ correction & Signal killed: $\sim 10^{-27}$~s,
  $10^{15}\times$ below detection &
  \textcolor{falsified}{\textbf{DEAD}} \\
W3i6 & Two-component decomposition: T5 resolved &
  Theoretically self-consistent; experimentally unchanged &
  \textcolor{falsified}{\textbf{DEAD}} (theory improved) \\
W3i7 & All corrections propagated; final assessment &
  \emph{This document} &
  \textcolor{falsified}{\textbf{FINAL}} \\
\bottomrule
\end{tabular}
\end{table}

%=============================================================================
\section{What Survives}
\label{sec:P1_survives}
%=============================================================================

\begin{tcolorbox}[colback=green!3!white, colframe=green!50!black,
  title=\textbf{P1: What Survives (Theory Only)}]

\begin{enumerate}[label=\textbf{S\arabic*.}, itemsep=4pt]

\item \textbf{The DFD refractive index prediction.}
  $n(r) = e^{\psi(r)}$, where $\psi$ is the DFD scalar field.
  In the two-component picture: $n = e^{\psigrav + \dpsiEM}$.
  This is a clean, testable prediction that survives all corrections.
  It predicts that regions of high EM energy density have a
  (vanishingly small) refractive index perturbation. [\textbf{VERIFIED}]

\item \textbf{The T5 resolution via two-component decomposition.}
  The EM-sourced $\dpsiEM$ legitimately modifies the spacetime metric
  (via $g_{00} = -e^{2\psi}$), resolving the apparent tension between
  EM-only sourcing and metric-coupling. P1's predictions are
  internally self-consistent within DFD. [\textbf{VERIFIED}]

\item \textbf{Non-reciprocal light travel time.}
  In a region with $\nabla\dpsiEM \neq 0$, light travel time
  is non-reciprocal: $\delta t_{A\to B} \neq \delta t_{B\to A}$.
  This is a distinctive DFD signature (GR gives reciprocal
  light travel in a static field). The non-reciprocity is:
  \begin{equation}
  \delta t_{\rm NR} = \frac{2}{c}\int_A^B \dpsiEM\,dl
  \sim \frac{2\,\dpsiEM\,L}{c} \sim \frac{2\times 10^{-22}\times 1}
  {3\ee{8}} \sim 7\ee{-31}~\text{s}.
  \end{equation}
  Undetectable, but the prediction survives as a theoretical
  consequence of DFD. [\textbf{VERIFIED, UNDETECTABLE}]

\item \textbf{Neutron star timing anomaly: 6.0~ps.}
  Near a neutron star ($\psigrav \sim 0.2$), the EM-sourced
  $\dpsiEM$ from the magnetosphere produces a timing perturbation
  of $\sim 6.0$~ps (W3i5-P1). This is enhanced by $e^{\psigrav}
  \approx 1.22$ (the gravitational--EM cross-term, Finding~3.6 of
  W3i6). Conditional on T5: if the two-component decomposition is
  correct, this prediction stands. Requires SKA Phase~2
  ($\sim 2030$s) for detection. [\textbf{CONDITIONAL, SURVIVING}]

\item \textbf{The drag-reduction effect.}
  A body moving through a region of modified $\psi$ experiences
  reduced gravitational drag (the $\psi$-modified geodesic deviates
  from the GR geodesic). The required $\dpsiEM > 10^{-10}$ for any
  measurable drag reduction requires $u_{\rm EM} > 10^{14}$~J/m$^3$
  sustained in macroscopic volumes. This is beyond any known or
  foreseeable technology by $>10^{10}\times$.
  [\textbf{THEORETICAL ONLY}]

\item \textbf{Gravitational cloaking.}
  Complete gravitational cloaking (making an object invisible to
  gravitational detection) requires $\dpsiEM \sim -\psigrav$ locally,
  which means $\dpsiEM \sim -GM/(c^2 r) \sim -10^{-9}$ for a 1-tonne
  object at 1~m. This requires $u_{\rm EM} \sim 10^{5}$~J/m$^3$ at
  $\lbone \sim 10^{-9}$ (at the SQMS bound). At first glance this
  seems achievable (SRF cavities reach $10^4$~J/m$^3$), but the
  actual requirement is:
  \begin{equation}
  \dpsiEM = (\lam-1)\,\frac{u_{\rm EM}}{c^2\,\rho_\psi}
  \sim 10^{-9}\times\frac{10^5}{9\ee{16}\times 10^{15}}
  \sim 10^{-36}.
  \end{equation}
  The density $\rho_\psi = \Meff/\Vcav$ enters, and the cloaking
  requires $\dpsiEM \sim 10^{-9}$, while the actual $\dpsiEM$
  from any achievable $u_{\rm EM}$ is $\sim 10^{-22}$. The gap
  is $10^{13}\times$. [\textbf{THEORETICAL ONLY}]
\end{enumerate}
\end{tcolorbox}

%=============================================================================
\section{What Is Dead}
\label{sec:P1_dead}
%=============================================================================

\begin{tcolorbox}[colback=red!5!white, colframe=red!60!black,
  title=\textbf{P1: What Is Dead (Experimentally Infeasible)}]

\begin{enumerate}[label=\textbf{D\arabic*.}, itemsep=4pt]

\item \textbf{The P1 timing experiment (4.0~ns signal).}
  The original P1 signal estimate was based on $\psi_{\rm cosmo}
  \approx 0.047$ (gravitational), which was applied to the EM-sourced
  perturbation. Under EM-only sourcing: $\dpsiEM^{\rm cosmo}
  = 9.6\ee{-6}$, and the actual $\dpsiEM$ from an SRF cavity is
  $\sim 10^{-22}$. Signal: $\delta t \sim 10^{-27}$~s. Required
  clock precision: $10^{-27}$~s (current best: $\sim 10^{-18}$~s,
  gap of $10^{9}\times$). [\textbf{DEAD, W3i5}]

\item \textbf{Laboratory drag reduction.}
  Any measurable drag effect requires $\dpsiEM > 10^{-10}$ in a
  macroscopic volume. Achievable $\dpsiEM \sim 10^{-22}$. Gap:
  $10^{12}\times$. [\textbf{DEAD, W3i5}]

\item \textbf{Laboratory gravitational cloaking.}
  Requires $\dpsiEM \sim 10^{-9}$. Achievable: $\sim 10^{-22}$.
  Gap: $10^{13}\times$. [\textbf{DEAD, W3i5}]

\item \textbf{Kelvin--Helmholtz suppression experiment.}
  The K-H suppression threshold requires a shear velocity modification
  that depends on $\nabla\dpsiEM$. With $\dpsiEM \sim 10^{-22}$ over
  $\sim 1$~m, $\nabla\dpsiEM \sim 10^{-22}$~m$^{-1}$. The K-H
  threshold modification is $\propto (\nabla\dpsiEM)^2 \sim 10^{-44}$.
  No conceivable experiment. [\textbf{DEAD, W3i4}]

\item \textbf{Any active P1 experiment.}
  All active P1 experiments require creating a macroscopic
  $\dpsiEM$ perturbation and detecting its effect on light
  propagation, mass motion, or fluid dynamics. The fundamental
  bottleneck is $|\dpsiEM|_{\rm lab} \sim 10^{-22}$
  (Theorem~4.7 of W3i6). No transduction chain can bridge
  $10^{12}$--$10^{15}$ orders of magnitude.
  [\textbf{DEAD, STRUCTURAL}]
\end{enumerate}
\end{tcolorbox}

%=============================================================================
\section{What Is Theoretical Only}
\label{sec:P1_theoretical}
%=============================================================================

\begin{tcolorbox}[colback=yellow!5!white, colframe=yellow!60!black,
  title=\textbf{P1: Theoretical Predictions (No Near-Term Test)}]

\begin{enumerate}[label=\textbf{T\arabic*.}, itemsep=4pt]

\item \textbf{The refractive index formula $n = e^\psi$.}
  This is a foundational DFD prediction. If DFD is confirmed by
  passive detection (SQMS, ESS, LIGO), the refractive formula is
  implicitly confirmed. No separate P1 experiment is needed to
  test it.

\item \textbf{Non-reciprocal light propagation in $\nabla\psi \neq 0$
  regions.}
  A genuine prediction of DFD that distinguishes it from GR.
  Detectable only near compact objects (neutron stars, black holes).
  Laboratory test: requires $\nabla\dpsiEM > 10^{-15}$~m$^{-1}$
  (achievable: $\sim 10^{-22}$~m$^{-1}$). Gap: $10^{7}\times$.

\item \textbf{Gravitational drag modification.}
  The DFD geodesic equation in a $\psi$-modified metric produces
  different orbital mechanics from GR. The difference is
  $\order{\dpsiEM^2}$ and undetectable in the solar system
  ($\dpsiEM^{\rm Sun} \sim 10^{-12}$ at Earth; correction
  $\sim 10^{-24}$ of GR effects).

\item \textbf{P1 as a technology demonstrator (far future).}
  If field-generation technology advances by $>10^{12}\times$
  (e.g., through breakthroughs in high-energy-density EM confinement,
  or if $\lbone$ is found to be much larger than the SQMS bound in
  a specific frequency range), P1-type experiments become possible.
  No timeline can be assigned. This is a ``centuries-scale''
  aspiration.
\end{enumerate}
\end{tcolorbox}

%=============================================================================
\section{P1 Final Verdict}
\label{sec:P1_verdict}
%=============================================================================

\begin{tcolorbox}[colback=blue!5!white, colframe=blue!60!black,
  title=\textbf{PATENT 1: DEFINITIVE FINAL ASSESSMENT (W3i7)}]

\textbf{Executive summary}: Patent 1 is theoretically self-consistent
within the DFD framework (after the W3i5 two-component T5 resolution)
but experimentally dead by $10^{12}$--$10^{15}$ orders of magnitude.
No foreseeable technology can close this gap. P1 retains value as:
\begin{itemize}[nosep]
\item A theoretical prediction of DFD that would be implicitly confirmed
  if DFD is validated by passive detection.
\item The motivation for the refractive-index formulation $n = e^\psi$
  that underpins the LIGO 6th-channel search (via the gradient coupling).
\item A far-future technology target if EM energy density or psi-field
  coupling exceeds current bounds by many orders of magnitude.
\end{itemize}

\textbf{Recommendation}: \textbf{Zero resources should be allocated
to P1-specific experimental work.} All P1-relevant physics is tested
more sensitively by the passive detection programme (SQMS Phase~I,
ESS Lund, LIGO 6th channel). P1 should be maintained as a theoretical
appendix to the DFD framework, not as an active research direction.

\textbf{Patent value}: P1's intellectual property value is in the
\emph{concept} of psi-field-mediated drag reduction and cloaking.
If DFD is confirmed and $\lbone$ is measured, P1 defines the
engineering target for future applications. The patent protection
is forward-looking, not near-term.

\begin{center}
\renewcommand{\arraystretch}{1.3}
\begin{tabular}{@{}lll@{}}
\toprule
\textbf{Aspect} & \textbf{Status} & \textbf{Confidence} \\
\midrule
Theoretical consistency & Self-consistent (T5 resolved) & High \\
Experimental feasibility & Dead ($10^{12}$--$10^{15}$ gap) & Definitive \\
Astrophysical test (NS timing) & Conditional, 6.0~ps & Medium \\
Implicit test via passive & Yes (if DFD confirmed) & High \\
Near-term resources & Zero recommended & Definitive \\
Patent value & Forward-looking IP & Medium \\
\bottomrule
\end{tabular}
\end{center}

\textbf{This assessment is FINAL barring a fundamental revision of
the DFD field equation, the discovery of a new coupling mechanism,
or a measurement of $\lbone > 10^{-5}$ that reopens the active
detection window.}
\end{tcolorbox}


%=============================================================================
%=============================================================================
\part{Combined W3i7 Findings, Corrections, and Open Questions}
\label{part:combined}
%=============================================================================
%=============================================================================

%=============================================================================
\section{Ranked Findings (W3i7)}
\label{sec:ranked_findings}
%=============================================================================

\begin{enumerate}[label=\textbf{W3i7-F\arabic*.}, itemsep=6pt]

\item \textbf{[RANK 1 --- CORRECTION] LIGO 6th-channel reach corrected
  from $6\ee{-19}$ to $\sim 3\ee{-14}$.}
  The common-mode rejection of scalar signals in a differential
  interferometer was not accounted for in W3i5/W3i6. The gradient
  coupling channel provides the dominant sensitivity, but at 5 orders
  of magnitude weaker than previously claimed. ESS Lund now leads
  the detection ranking.

\item \textbf{[RANK 2 --- NEW] Complete LIGO systematic budget derived.}
  Dominant systematics: Newtonian noise (indistinguishable from psi
  below 20~Hz), Schumann resonances (correlated across detectors,
  degenerate with solar psi), laser frequency noise (common-mode
  mimic). Five null tests defined: polarisation, EM--GW correlation,
  dispersion, anisotropy, diurnal modulation.

\item \textbf{[RANK 3 --- NEW] SQMS Phase~I experimental protocol
  specified.}
  $2\times 2$ cavity arrangement with intra-pair/inter-pair
  correlation diagnostics. 6-month baseline characterisation followed
  by 12-month science data taking. Four-criterion detection threshold.
  Multi-frequency $Q_0$ measurement identified as most powerful
  single diagnostic.

\item \textbf{[RANK 4 --- VERIFIED] Existing SQMS data provide
  only trivial $\Qpsi$ constraints.}
  The $\lbone/\Qpsi$ degeneracy cannot be broken with passive
  measurements alone. Multi-frequency and field-dependent $Q_0$
  measurements provide model-dependent handles but not clean
  constraints. $\Qpsi$ measurement requires prior detection.

\item \textbf{[RANK 5 --- NEW] Five existing datasets identified
  for reanalysis.}
  SQMS multi-frequency $Q_0$, DESY TESLA database, JLab C100
  correlations, LIGO cavity frequency, PTB sapphire oscillators.
  Multi-frequency residual ratio is the most promising.

\item \textbf{[RANK 6 --- FINAL] P1 definitive assessment complete.}
  Theoretically self-consistent (T5 resolved). Experimentally dead
  by $10^{12}$--$10^{15}\times$. NS timing anomaly (6.0~ps) is only
  surviving near-term testable prediction (SKA Phase~2, 2030s).
  Zero resources recommended for P1-specific experiments.

\item \textbf{[RANK 7 --- CORRECTED] Updated detection ranking.}
  ESS Lund optimistic ($10^{-14}$) now \#1; LIGO ($3\ee{-14}$) now
  \#2; ESS realistic ($10^{-12}$) \#3; SQMS current ($1.56\ee{-9}$) \#4.
  Passive--active gap reduced from 14 to 9 orders of magnitude but
  remains fundamental.
\end{enumerate}

%=============================================================================
\section{W3i7 Corrections to Previous Iterations}
\label{sec:w3i7_corrections}
%=============================================================================

\begin{table}[H]
\centering
\caption{W3i7 corrections relative to W3i6.}
\label{tab:w3i7_corrections}
\renewcommand{\arraystretch}{1.25}
\begin{tabular}{@{}p{4cm}p{3cm}p{3cm}l@{}}
\toprule
\textbf{Item} & \textbf{W3i6 Value} & \textbf{W3i7 Value} &
  \textbf{Status} \\
\midrule
LIGO 6th-channel reach & $6\ee{-19}$ & $\sim 3\ee{-14}$ &
  \textcolor{falsified}{\textbf{CORRECTED $10^5\times$}} \\
Detection ranking \#1 & LIGO & ESS Lund (optimistic) &
  \textcolor{conjcolor}{\textbf{RERANKED}} \\
Passive--active gap & $10^{14}\times$ & $10^{9}\times$ &
  \textcolor{conjcolor}{\textbf{CORRECTED}} \\
$\Qpsi$ constraint from data & ``Not proposed'' & ``Trivial, cannot
  break degeneracy'' &
  \textcolor{falsified}{\textbf{ASSESSED}} \\
P1 status & Dead (experimental) & \textbf{FINAL: dead experimental,
  self-consistent theory} &
  \textcolor{falsified}{\textbf{CLOSED}} \\
SQMS Phase~I protocol & ``4-cavity, \$2--4M'' & Detailed protocol
  with measurement plan, detection criteria &
  \textcolor{strongcolor}{\textbf{ELABORATED}} \\
LIGO systematics & ``Extreme systematic risk'' & Complete budget with
  5 null tests &
  \textcolor{strongcolor}{\textbf{NEW}} \\
\bottomrule
\end{tabular}
\end{table}

%=============================================================================
\section{Impact on the Five-Step Programme}
\label{sec:five_step_update}
%=============================================================================

The W3i6 five-step programme is modified:

\begin{tcolorbox}[colback=blue!5!white, colframe=blue!50!black,
  title=\textbf{W3i7 Updated Five-Step Programme}]

\textbf{Step 1: SQMS Phase~I (2027--2028)} --- \textbf{PROMOTED to
co-equal with LIGO.}
  With LIGO reach corrected to $\sim 10^{-14}$ (comparable to
  ESS optimistic), the SQMS 4-cavity array becomes more strategically
  important: it is the \emph{only} experiment that can be built
  entirely within a single institution's control. Protocol specified
  in Part~\ref{part:SQMS}.

\textbf{Step 2: LIGO 6th Channel (2027--2029)} --- \textbf{DEMOTED
from \#1 to parallel track.}
  Still valuable at $\sim 10^{-14}$, but no longer the ``silver
  bullet.'' Systematic budget (Part~\ref{part:LIGO}) shows significant
  obstacles. Proceed as low-cost data analysis effort.

\textbf{Step 3: ESS Lund (Q2 2027)} --- \textbf{PROMOTED to highest
  priority external collaboration.}
  Now the single most sensitive projected experiment. The lock-in
  measurement at 14~Hz beam frequency bypasses most of the LIGO
  systematics. Institutional buy-in is the rate-limiting step.

\textbf{Step 4: Indirect $\Qpsi$ (2027--2028)} --- \textbf{ASSESSED as
  unable to constrain $\Qpsi$ without prior detection.}
  Multi-frequency $Q_0$ ratio at SQMS is the best diagnostic but cannot
  break the $\lbone/\Qpsi$ degeneracy alone. Reanalysis of 5 existing
  datasets recommended (Finding~\ref{find:datasets}).

\textbf{Step 5: MEMS+P11 (2028+, conditional)} --- \textbf{UNCHANGED.}
  Still conditional on Steps 1--4 detecting $\lbone \neq 0$ and
  $\Qpsi > 10^{10}$.
\end{tcolorbox}

%=============================================================================
\section{Open Questions for W3i8}
\label{sec:open_questions}
%=============================================================================

\begin{enumerate}[label=\textbf{OQ-\arabic*.}, itemsep=4pt]

\item \textbf{ESS Lund LLRF systematic floor}: What is the achievable
  $\delta Q/Q$ precision at 14~Hz from the ESS digital LLRF system?
  Can it reach $10^{-6}$ or is $10^{-4}$ the practical limit?
  (Requires P3/P10 agent input.)

\item \textbf{LIGO scalar mode search}: Has the LIGO collaboration
  published any scalar-mode search that can be reinterpreted for DFD?
  The existing ``breathing mode'' searches for Brans-Dicke scalar GW
  may directly constrain $\dpsiEM$. (Requires P6 agent input.)

\item \textbf{Multi-frequency $\Heff$ calculation}: The $\Heff$ values
  used in Finding~\ref{find:multi_freq} are estimates. A full EM
  field simulation of TESLA cavity modes with psi-mode overlap
  integrals is needed for the Phase~I analysis pipeline.

\item \textbf{Newtonian noise subtraction for scalar modes}: The
  existing NN subtraction literature focuses on tensor GW.
  Is the transfer function from seismic to scalar strain different?
  Can the same seismometer arrays be used?

\item \textbf{Cross-reference with P14 (shadow)}: The $+4.6\%$ shadow
  prediction operates in the $\psigrav$ sector. If EHT confirms the
  shadow excess, does this provide any constraint on $\lbone$ or
  $\Qpsi$ indirectly? (Through the two-component self-consistency.)

\item \textbf{P1 patent strategy}: With P1 experimentally closed,
  what is the IP strategy? Maintain broad claims for future technology,
  or narrow to specific theoretical predictions? (Requires legal input.)

\end{enumerate}

%=============================================================================
\section*{Literature References}
%=============================================================================

\begin{thebibliography}{99}

\bibitem{W3i6_P1P2}
DFD Research Programme, Agent~1/2,
``W3i6: Post-Paradigm-Shift Signal Ranking, Two-Component Physics for P1,
Passive Detection Strategy for P2,''
\textit{DFD Research Output}, W3i6\_P1\_P2\_update.tex, March 2026.

\bibitem{W3i6_P8P11}
DFD Research Programme, Agent~8/11,
``W3i6: Two-Component Impact on P8, Definitive Detection Roadmap for P11,''
\textit{DFD Research Output}, W3i6\_P8\_P11\_update.tex, March 2026.

\bibitem{W3i6_MatSci}
DFD Research Programme, Materials Science Agent,
``W3i6: Passive Material FOM, MEMS Optimisation, Updated Roadmap,''
\textit{DFD Research Output}, W3i6\_MatSci\_update.tex, March 2026.

\bibitem{W3i5_P2}
DFD Research Programme, Agent~2,
``W3i5: Passive Superiority Theorem, Volume Cancellation Theorem,
$Q_\psi$ Threshold,''
\textit{DFD Research Output}, W3i5\_P2\_update.tex, March 2026.

\bibitem{W3i5_P1}
DFD Research Programme, Agent~1,
``W3i5: $M_\psi$ Correction, T1/T5 Impact, P1 Signal Killed,''
\textit{DFD Research Output}, W3i5\_P1\_update.tex, March 2026.

\bibitem{Romanenko2020}
A.~Romanenko, A.~Grassellino, et al.,
``Ultra-high quality factors in superconducting niobium cavities in
ambient magnetic fields,''
\textit{Phys.\ Rev.\ Applied} \textbf{13}, 034032 (2020).

\bibitem{LIGO_scalar}
B.~P.~Abbott et al.\ (LIGO/Virgo),
``Tests of general relativity with binary black holes from the second
LIGO-Virgo gravitational-wave transient catalog,''
\textit{Phys.\ Rev.\ D} \textbf{103}, 122002 (2021).

\bibitem{Newtonian_noise}
J.~Harms,
``Terrestrial gravity fluctuations,''
\textit{Living Rev.\ Relativ.} \textbf{22}, 6 (2019).

\bibitem{Schumann}
E.~Thrane, N.~Christensen, and R.~Schofield,
``Correlated magnetic noise in global networks of gravitational-wave
detectors: observations and implications,''
\textit{Phys.\ Rev.\ D} \textbf{87}, 123009 (2013).

\bibitem{SQMS_Qfactor}
A.~Grassellino et al.,
``Nitrogen and argon doping of SRF cavities for high Q factors,''
\textit{Supercond.\ Sci.\ Technol.} \textbf{30}, 094004 (2017).

\end{thebibliography}

\end{document}
